\documentclass[conference]{IEEEtran}
\usepackage[utf8]{inputenc}
\usepackage[T1]{fontenc}
\usepackage{cite}
\usepackage{amsmath,amssymb}
\usepackage{graphicx}
\usepackage{booktabs}
\usepackage{url}

\title{Preregistered Paired Evaluation of a Deterministic Verifier Pipeline for LLM‑Generated Network Configuration Changes --- Non‑informative Result under Shared‑Generation Semantics}

\author{
\IEEEauthorblockN{Autonomous AI Researcher}
\IEEEauthorblockA{CNSM 2026 Agentic AI Researcher Track}
}

\begin{document}

\maketitle

\begin{abstract}
We report a fully preregistered paired experiment (cand-2026-01-prereg-v1) that tested whether a deterministic verifier pipeline (generate \ensuremath{\rightarrow} deterministic\_netops\_validator\_v5 \ensuremath{\rightarrow} optional repair) reduces the rate of verifier‑detected unsafe intent\ensuremath{\rightarrow}configuration proposals produced by a hosted LLM on synthetic NetOps tasks. Design: N=200 preregistered unique intent\ensuremath{\rightarrow}configuration tasks in a paired design comparing (a) baseline LLM‑only and (b) guarded pipeline (gpt-5-mini + deterministic\_netops\_validator\_v5). Execution used shared\_initial\_candidate generation semantics (one LLM call per task) producing model\_calls\_used = 200 (preregistered independent per‑arm generation would have used up to 400). Observed (adversarial cluster): baseline SAFE = 200/200; guarded SAFE = 200/200 (paired contingency n11 = 200, n10 = 0, n01 = 0). Estimated paired SAFE‑rate difference = 0.0; exact McNemar two‑sided p = 1.0; paired bootstrap 95\% CI = [0.0, 0.0] (resamples = 10000, seed = 1729). Because identical per‑pair generations were produced and the observed baseline unsafe rate was 0\%, the preregistered causal hypothesis (that the deterministic verifier reduces unsafe proposals) was not testable in this execution. We publish the complete preregistered run and artifacts, provide an artifact‑grounded diagnosis of testability failure, and give concrete, archive‑supported recommendations (independent per‑arm generation budgeting, adversarial injection calibration, and exploratory ROC reserves) for informative repeats. All execution and analysis artifacts cited here are archived in the supplied bundle.
\end{abstract}

\section{Introduction}
Automating human‑intent \ensuremath{\rightarrow} device‑configuration translation with large language models (LLMs) is an active NetOps research thrust because it can reduce human effort and speed maintenance tasks. Yet incorrectly specified or misordered changes can cause outages, policy violations, or security exposures. A commonly proposed mitigation is tool‑augmentation: couple an LLM generator to deterministic validators, sandboxed simulators, or constrained executors in a generate--validate--repair pipeline so that proposed changes are checked before execution.

We preregistered and executed a confirmatory paired experiment (cand-2026-01-prereg-v1) to quantify whether deterministic verifier augmentation causally reduces the frequency of verifier‑detected unsafe proposals from a hosted LLM under both benign and adversarial prompt clusters. The preregistered estimand was the paired SAFE‑rate difference (guarded minus baseline), tested by an exact McNemar two‑sided test with a paired bootstrap CI (10000 resamples, seed = 1729). The design sampled N=200 tasks using a deterministic generator and a paired comparison across two pipelines. Two generation semantics were permitted in the preregistration: independent per‑arm generation (one LLM call per arm per task) or shared\_initial\_candidate (one shared LLM call per task scored by both arms). The executed run used shared\_initial\_candidate to conserve model‑call budget.

Key outcome: under the executed shared semantics the sampled LLM outputs were scored SAFE by the deterministic validator in every confirmatory episode (baseline SAFE = 200/200; guarded SAFE = 200/200). This concordant ceiling outcome (n11 = 200; n10 = n01 = n00 = 0) yields an estimate of zero paired difference and an exact McNemar p = 1.0. Because the observed baseline unsafe rate was 0\% and the same candidate was used in both arms, the preregistered causal hypothesis could not be meaningfully evaluated. The manuscript therefore focuses on three contributions grounded in archived artifacts: (1) a complete, deterministic preregistered run published as reproducible artifacts; (2) a precise, artifact‑grounded diagnosis of why this execution was non‑informative for the confirmatory estimand; and (3) concrete, artifact‑supported design recommendations that would enable informative repeats under the same preregistration framework.

\section{Related Work}
Recent surveys and prototypes document strong interest in applying LLMs to NetOps tasks---intent translation, configuration synthesis, diagnosis, and limited remediation---but emphasize that demonstrated gains are largely in prototype or curated‑case settings rather than in production deployments [openalex:W7161354925; openalex:W7170714602; openalex:W4411109869]. Several lines of prior work are directly relevant to our design choices and to the generate--validate--repair architecture we evaluated. First, multiple studies and proposals show that augmenting LLMs with deterministic tools (verifiers, simulators, or formal checkers) reduces certain classes of constraint violations on curated tasks: generate--validate--refine or divide--verify--refine pipelines measurably improve adherence to structured instructions when the verifier faithfully captures the task constraints [openalex:W4412888330; openalex:W7161354925]. These works, however, also report that improvement magnitude varies with verifier fidelity and the nature of the constraints being checked.

Second, research on guardrails, tool schemas, and MCP‑style governance highlights practical trade‑offs encountered by tool‑augmented agents: stricter schema enforcement reduces some unsafe outputs but can increase false refusals or operator workload unless tuned carefully [TechRxiv:177006109; TechRxiv:177004277]. Red‑teaming and guardrail‑distillation literature from adjacent domains shows that adversarial prompts can often be engineered to bypass naive guardrails, motivating replayable, sandboxed evaluations and explicit ROC‑style diagnostics for verifiers [openalex:W7172000952; TechRxiv:177006109].

Third, evaluation critiques argue that static benchmarks are insufficient for agentic NetOps: meaningful assessment requires workflow‑centred, deterministic replay (sandbox) environments, archived execution traces, and paired causal designs that can separate generator variability from validator effects [openalex:W7161354925; doi:10.2139/ssrn.7268438]. These methodological prescriptions informed our preregistration (paired design, deterministic task generator, and archived validator traces) because they make it possible to ascribe observed outcome differences to pipeline interventions rather than to uncontrolled environmental variability.

Finally, prototype and architecture surveys stress operational integration concerns---human‑in‑the‑loop gates, auditability, multi‑tool orchestration, and simulator fidelity---all of which constrain what a verifier can safely assert in practice and how empirical evaluation results should be interpreted for operational readiness [TechRxiv:177004277; openalex:W4406549441]. In sum, prior work supports the conceptual value of verifier augmentation but also warns that empirical gains are task‑ and verifier‑dependent and that evaluation semantics (for example, shared vs independent generation and verifier completeness) strongly affect the interpretability of paired causal contrasts. Our preregistered design attempted to operationalize these recommendations; the archived run and the diagnosis below show how a specific execution semantics (shared\_initial\_candidate) interacted with task difficulty and verifier outputs to produce a non‑informative confirmatory outcome.

\section{Methodology}
Preregistration and primary estimand. The experiment was preregistered as cand‑2026‑01‑prereg‑v1. The primary estimand is the paired binary difference paired\_success\_rate\_difference\_guarded\_minus\_baseline where each arm contributes a binary indicator per task (1 = SAFE, 0 = UNSAFE) derived from a deterministic sandbox validator. The confirmatory test is an exact two‑sided McNemar implemented by the paired\_binary\_analysis\_v1 executor; paired bootstrap percentile confidence intervals use 10,000 resamples with seed = 1729. The preregistration specified a paired binary design with N=200 unique tasks (400 paired episodes across two arms under independent per‑arm draws), power calculations targeting an \ensuremath{\approx}7 percentage‑point paired reduction, and predeclared sensitivity and multiplicity procedures.

Task generation, scope, and deterministic seeding. A deterministic task generator (deterministic\_netops\_task\_generator\_v5) produced 200 synthetic intent\ensuremath{\rightarrow}configuration tasks constrained to allowed fields \{admin, mtu, vlan\} within single‑AS multi‑device topologies. Each task manifest includes initial\_state, expected\_final\_state, an explicit required\_sequence (workflow policy), and difficulty metadata; generation was seeded deterministically (per‑task generation\_seed 1..200 as recorded in execution/task\_manifest.jsonl) so every task can be reconstructed exactly. Tasks were classified into benign and adversarial clusters using a preregistered adversarial template set designed to increase difficulty by obfuscation, contradictory constraints, and privilege‑escalation phrasing; the adversarial cluster formed the confirmatory challenge set alongside benign items per the preregistration.

Pipelines compared and validator semantics. Two pipelines were contrasted: (1) baseline (LLM‑only): accept the generator output and evaluate post hoc with deterministic\_netops\_validator\_v5; (2) guarded (LLM + deterministic validator): validate the generator output using deterministic\_netops\_validator\_v5 in a sandboxed replay and (optionally) invoke an automatic repair step once. For confirmatory episodes in this run the repair pathway was not exercised (repair\_provider\_trace fields are null in execution/scoring/*.json), so the guarded arm operationally reflects deterministic validation only. The validator applies the preregistered scoring\_rule full\_intent\_constraint\_validation and emits deterministic traces (validation\_before, validation\_after, violation lists, and a boolean valid flag) that produce the binary SAFE/UNSAFE outcome used by analysis. Representative validator traces are archived under execution/scoring/* and execution/responses/* and show the validator\_id and validator\_version recorded per episode.

Generation semantics, model budgeting, and execution accounting. The preregistration allowed two generation semantics: independent per‑arm generation (two independent LLM calls per task, one for each arm) or shared\_initial\_candidate (one shared LLM call per task whose candidate is scored by both arms). To conserve hosted model budget the executed run used shared\_initial\_candidate semantics; this is recorded in the archived model configuration (execution/model\_configuration.json: independent\_condition\_generation = false) and in the execution manifest (execution/execution\_manifest.json: model\_calls\_used = 200). The preregistered independent per‑arm design would have used up to 400 model calls (two per task) under the same N=200 confirmatory sample; that budget differential is therefore the precise additional model‑call requirement to enable arm‑level stochastic contrasts under the preregistered plan. The execution manifest and model configuration are archived and document these accounting facts.

Scoring, missingness, exclusions, and deterministic reconciliation. Deterministic\_netops\_validator\_v5 produced a deterministic valid flag per candidate; confirmatory binary labels were derived directly from validator.valid. The preregistration specified one automatic retry for model call failures, duplicate‑prompt detection (prompt hashes lead to automatic exclusion+replacement), verifier nondeterminism checks (replays), and a stopping rule that aborts or downgrades if technical failures exceed 10\%. In the archived confirmatory run, deterministic reconciliation artifacts (analysis/deterministic\_reconciliation.json and execution/execution\_manifest.json) report completed\_episode\_count = 400, failed\_episode\_count = 0, complete\_pairs = 200, and that all deterministic consistency checks passed; analysis/contamination\_summary.csv and analysis/missingness\_summary.csv document zero exclusions or technical failures. Those artifacts confirm that the run executed under the shared generation semantics without technical interruptions and that no contaminated or nondeterministic tasks were present in the confirmatory set.

Analysis pipeline and confirmatory computation. The analysis executor paired\_binary\_analysis\_v1 consumed the per‑episode binary outcomes and produced the 2\ensuremath{\times}2 contingency (analysis/paired\_contingency\_table.csv: n11,n10,n01,n00), exact McNemar two‑sided p‑value, and paired bootstrap CI (analysis/analysis\_log.jsonl). For this execution the contingency was n11 = 200, n10 = 0, n01 = 0, n00 = 0, yielding an estimated paired SAFE‑rate difference = 0.0, exact McNemar p = 1.0, and paired bootstrap 95\% CI = [0.0, 0.0]. Representative per‑task scoring JSONs and shared initial responses are archived under execution/scoring/ and execution/responses/ (e.g., task-000001‑*.json and task-000001‑shared-initial.txt) and show identical shared\_initial\_candidate content and validator traces with violation\_count = 0. The preregistered exploratory reserve (up to 50 tasks to estimate verifier ROC/refusal trade‑offs) was not consumed prior to confirmatory execution; consequently no ROC or refusal estimates appear in the confirmatory artifacts.

Design consequences and protocol fidelity. Because the executed semantics produced a single shared candidate per task that both arms scored deterministically and because the validator labeled every sampled candidate SAFE, the paired design produced no discordant pairs and therefore no testable paired effect under McNemar logic. This outcome is an execution‑semantic and budgeting consequence that follows directly from archived configuration and execution artifacts (execution/model\_configuration.json and execution/execution\_manifest.json) rather than from properties of the validator algorithm. All files referenced above are included in the archived bundle so readers can replay generation, scoring, and analysis deterministically to confirm the chain of events described here.

\section{Results}
Execution accounting and provider audit. The execution manifest reports completed\_episode\_count = 400 (200 paired episodes), failed\_episode\_count = 0, and model\_calls\_used = 200, consistent with the executed shared\_initial\_candidate semantics. Provider‑level auditing recorded hosted\_model\_call\_event\_count = 200 and cache\_miss\_count = 200; cross‑condition cache reuse was not observed. No technical failures, exclusions, or nondeterministic verifier behaviors were recorded in the confirmatory set, and deterministic reconciliation checks passed, confirming that the archived episode and scoring artifacts correspond to a single, reproducible run.

Primary confirmatory outcomes and contingency. The deterministic validator produced SAFE labels for every confirmatory episode: baseline SAFE successes = 200/200 and guarded SAFE successes = 200/200. The resulting paired contingency counts are n11 = 200 (SAFE in both arms), n10 = 0 (SAFE guarded only), n01 = 0 (SAFE baseline only), and n00 = 0 (UNSAFE in both). From these counts the paired SAFE‑rate difference (guarded - baseline) equals 0.0. The preregistered exact McNemar two‑sided test computed p = 1.0; the paired bootstrap percentile 95\% confidence interval is degenerate at [0.0, 0.0] (resamples = 10000; bootstrap\_seed = 1729). These numerical values are direct outputs of paired\_binary\_analysis\_v1 and are archived in the analysis artifacts (condition\_summary and paired\_contingency\_table files).

Representative diagnostic traces. Representative scoring JSONs and validator traces illustrate why the contingency is degenerate. For multiple sampled tasks the shared\_initial\_candidate content (the raw generated configuration) is identical across baseline and guarded episodes; validator traces show validation\_before and validation\_after with violation\_count = 0 and valid = true in every case. The repair\_provider\_trace entries are null in all confirmatory scoring JSONs, confirming that no automatic repair was invoked and that the guarded arm performed only validation. Shared response files for representative tasks contain the canonical generated configurations used by both arms, reinforcing that no arm‑specific generation stochasticity existed under the executed semantics.

Missingness, contamination, and exploratory reserves. Analysis/missingness\_summary documents complete\_pairs = 200 with zero failed or unscorable episodes; analysis/contamination\_summary shows no flagged episodes. The preregistration reserved up to 50 exploratory tasks to estimate verifier ROC and refusal trade‑offs; that exploratory reserve was not consumed prior to confirmatory execution, so no empirical ROC or refusal‑rate estimates appear in the confirmatory analysis artifacts.

Statistical interpretation and inferential consequence. The observed ceiling outcome (n11 = 200 with no discordant pairs) renders the preregistered inferential machinery non‑informative for the intended causal estimand: McNemar's test relies on discordant pairs (n10 and n01) to evaluate paired differences, and their absence produces a p‑value of 1.0 and a degenerate CI at zero. This degeneracy does not demonstrate verifier effectiveness; rather, it indicates that under the executed shared generation semantics and for the sampled model outputs, there was no arm‑level variation to attribute to the presence or absence of the deterministic validator. The preregistered power calculation had assumed a nonzero baseline unsafe rate (\ensuremath{\approx}10--15\%) and discordant proportions; those assumptions were not met in the executed run because budgeted generation semantics constrained per‑arm variability.

Archival fidelity and reproducibility pointers. All numerical summaries and the contingency table reported above are identical to the archived analysis artifacts (analysis/condition\_summary.csv and analysis/paired\_contingency\_table.csv). Representative per‑task scoring JSONs and shared response files show the identical generated candidate for baseline and guarded episodes and document validator decisions and violation lists. The execution and analysis logs capture the exact bootstrap parameters (resamples = 10000, seed = 1729) and the exact statistical outputs. Because the confirmatory run contains no technical failures or exclusions, the archived artifacts permit deterministic replay of the identical non‑informative outcome; investigators seeking an informative repeat should consult these artifacts before changing preregistered design choices.

\section{Discussion}
Artifact‑grounded diagnosis of the testability (ceiling) failure. Three concrete archived facts explain why the confirmatory hypothesis was untestable in this execution: (1) the run executed shared\_initial\_candidate semantics (one LLM call per task shared across arms) to conserve budget; (2) the sampled gpt-5-mini outputs were labeled SAFE by deterministic\_netops\_validator\_v5 for every confirmatory task (baseline SAFE = 200/200); and (3) no repair calls were applied in confirmatory episodes, so the guarded arm could not introduce arm‑only divergences. Together these facts yield a degenerate paired contingency with no discordant pairs, which makes McNemar inference vacuous (p = 1.0; CI degenerate at 0.0).

Interpretation relative to the preregistered power plan. The preregistration's power calculation assumed a nonzero baseline unsafe rate (\ensuremath{\approx}10--15\%) and discordant proportions that would render a paired reduction (\textasciitilde{}7 percentage points) detectable with N\ensuremath{\approx}174. The executed shared semantics halved model calls (200 vs planned 400 for independent per‑arm draws) and removed arm stochasticity; hence the empirical baseline unsafe rate observed here (0\%) contradicts preregistered power assumptions and collapses inferential sensitivity. This is an execution semantic and budgeting failure, not evidence about verifier efficacy.

Concrete, artifact‑supported recommendations for informative repeats. All recommendations below follow from archived planning and execution artifacts:
- Use independent per‑arm generation to enable arm‑level stochastic contrasts when the estimand is the effect of adding a deterministic validator to generation. For N=200 this requires \ensuremath{\approx}200 extra model calls (executed model\_calls\_used = 200 vs preregistered maximum\_model\_calls = 400). If repairs are not applied, independent draws are necessary to create discordant pairs.
- Reserve and use the preregistered exploratory budget (up to 50 tasks) to tune adversarial prompt transforms until a nonzero baseline unsafe rate consistent with power assumptions is observed. Archive ROC/reserve outputs prior to confirmatory execution.
- If model‑call budgets are constrained, prioritize independent per‑arm draws over additional model varieties because paired causal contrasts require per‑arm variability more than additional cross‑model comparisons within the same budget.
- When repairs are allowed in the guarded arm, instrument and archive repair traces and report how often repairs convert UNSAFE \ensuremath{\rightarrow} SAFE and whether repairs introduce new violations; include repair\_provider\_trace artifacts in analysis. In our confirmatory run repair traces were null and therefore not a source of discordance.

Threats to validity and generalizability. Internal validity: results depend critically on generation semantics and model budgeting; the executed shared semantics created the ceiling failure. Construct validity: deterministic\_netops\_validator\_v5 encodes a formalized full\_intent\_constraint\_validation for our synthetic task family, but its completeness relative to all real‑world NetOps failure modes is not established. External validity: tasks are synthetic single‑AS topologies limited to admin/mtu/vlan changes and the evaluation used a single declared model (gpt-5-mini); results do not generalize to multi‑vendor production networks or other LLMs. Conclusion validity: the observed degenerate contingency prevents inferential claims about verifier efficacy in this execution; the statistical machinery produced p = 1.0 and a degenerate CI because there were no discordant observations.

\section{Limitations}
\begin{itemize}
\item Synthetic tasks and scope: tasks were deterministic synthetic topologies produced by deterministic\_netops\_task\_generator\_v5 and limited to single‑AS, multi‑device setups; results may not generalize to production, multi‑AS, or vendor‑specific features.
\item Generation semantics: the executed run used shared\_initial\_candidate (independent\_condition\_generation = false), producing identical per‑pair generations and creating a ceiling effect when shared candidates were valid.
\item Adversarial strength: the preregistered adversarial templates used in this run did not elicit unsafe outputs from gpt‑5‑mini on the sampled tasks; stronger adversarial cases or explicit injected unsafe examples are required to measure verifier effect.
\item Single‑model scope: only the declared model (gpt‑5‑mini) was evaluated; no multi‑model generality is claimed.
\item Verifier completeness: deterministic\_netops\_validator\_v5 ran deterministically in synthetic sandboxed states; external validation of validator completeness against all real‑world NetOps failure modes is not provided.
\item Operational external validity: no production deployment, operator‑in‑the‑loop workload measurement, nor multi‑vendor field trials were performed; operator‑cost trade‑offs remain unquantified at scale.
\end{itemize}

\section{Conclusion}
We executed a fully archived preregistered paired experiment comparing gpt-5-mini baseline generation and a deterministic verifier‑augmented pipeline on 200 synthetic NetOps tasks. Under the executed shared\_initial\_candidate semantics and for the declared model, both arms produced zero verifier‑detected unsafe proposals (paired difference = 0.0; exact McNemar p = 1.0; 95\% CI = [0.0, 0.0]). Because identical per‑pair generations were used and because the observed baseline unsafe rate was 0\%, the preregistered causal hypothesis was not testable in this run. The scientific contribution is therefore methodological and reproducibility‑centred: (1) a complete deterministic preregistered run with archived artifacts that permit exact replay; (2) an artifact‑grounded diagnosis of why this execution was non‑informative for verifier efficacy; and (3) concrete, archive‑supported recommendations (independent per‑arm generation budgeting, adversarial calibration, and exploratory ROC estimation) for informative repeats. The archived execution and analysis artifacts cited throughout (execution/* and analysis/*) permit deterministic replication of the diagnosis and the run outputs and should be consulted prior to attempting repeats or production extrapolation.

\section*{Disclosure Statement}
Disclosure Statement (CNSM Agentic AI Researcher track): Declared LLM: gpt-5-mini (execution recorded in execution/model\_configuration.json). Orchestration/framework: hosted\_netops\_gvr\_v1 adapter and deterministic experiment harness (execution/execution\_manifest.json). Exact initial master prompt: archived at provenance/master\_prompt.txt with immutable SHA‑256 = 1872df1e1805d2d96940456ca016bd665d1d5196add77f5acdf1582bb39b15ba. Topic constraints and autonomy boundary: the human Research Director selected the broad topic family (Generative AI and LLMs for NetOps) prior to the autonomy lock; after the lock the autonomous researcher performed literature review, preregistration (cand-2026-01-prereg-v1), experiment design, code generation, execution, analysis, manuscript drafting, autonomous peer review, and revisions. Preregistration and execution: the preregistration was frozen before execution; key preregistered elements (estimand, paired design N=200, task generator, allowed generation semantics, scoring rules, and stopping rules) are recorded in cand-2026-01-prereg-v1 and archived in the manuscript bundle. Generation semantics and model‑call accounting (executed): independent\_condition\_generation = false (shared\_initial\_candidate semantics) producing a single initial LLM candidate per task shared across arms; model\_calls\_used = 200 while the preregistered maximum\_model\_calls allocated for independent per‑arm generation = 400 (execution/execution\_manifest.json and execution/model\_configuration.json). Validator and scoring: deterministic\_netops\_validator\_v5 performed deterministic sandboxed validation and encoded outcomes as binary labels (1 = SAFE, 0 = UNSAFE) under scoring\_rule = full\_intent\_constraint\_validation; archived scoring JSONs (execution/scoring/*.json) contain validator traces showing validation\_before/after and violation lists. Analysis executor: paired\_binary\_analysis\_v1 computed the paired contingency, exact McNemar two‑sided test, and paired bootstrap CI (resamples = 10000, seed = 1729); analysis artifacts include analysis/paired\_contingency\_table.csv and analysis/condition\_summary.csv. Deviations and restarts: no post‑preregistration design changes were executed; the observed model call count reflects the executed shared\_initial\_candidate semantics. Reproducibility pointers (concise): execution/raw\_results.jsonl, execution/task\_manifest.jsonl, execution/model\_configuration.json, execution/execution\_manifest.json, analysis/paired\_contingency\_table.csv, analysis/condition\_summary.csv, and analysis/paired\_difference\_figure.svg are archived in the supplied bundle and permit deterministic replays. Human editing: no manual human edits to technical content, analyses, or manuscript text occurred after the autonomy lock. Pre‑lock development: infrastructure rehearsals occurred prior to the lock but were excluded from final empirical evidence unless reproduced under the post‑lock execution.

\begin{thebibliography}{99}
\bibitem{ref1} https://arxiv.org/abs/2605.12729
\bibitem{ref2} https://doi.org/10.2139/ssrn.7132138
\bibitem{ref3} https://doi.org/10.36227/techrxiv.177004277.71795055/v1
\bibitem{ref4} https://doi.org/10.36227/techrxiv.177006109.97957916/v1
\bibitem{ref5} https://doi.org/10.1109/ojcoms.2026.3680457
\bibitem{ref6} Xianren Zhang, Xianfeng Tang, Hui Liu, Zongyu Wu, Qi He, Dongwon Lee, Suhang Wang, "Divide-Verify-Refine: Can LLMs Self-align with Complex Instructions?", 2025, 10.18653/v1/2025.findings-acl.709
\bibitem{ref7} Muhammad Bilal, Jon Crowcroft, Ruizhi Wang, Xiaolong Xu, Schahram Dustdar, "Large Language Models for Agentic NetOps and AIOps: Architectures, Evaluation, and Safety", 2026
\end{thebibliography}

\end{document}
